%% GABARIT D'ARTICLE POUR LA REVUE CANADIENNE DE STATISTIQUE
%%
%% Consultez les directives pour les auteurs de la RCS pour la
%% rédaction générale et la documentation de la classe cjs-rcs-article
%% pour une description détaillée de la classe.
%%
%% [N'hésitez pas à supprimer les commentaires après les avoir lus.]
%%
%% Déclaration de la classe pour un article en français. Vous devez
%% soumettre votre manuscript avec l'option 'review'. Supprimez
%% l'option seulement une fois que votre article aura été accepté pour
%% publication.
\documentclass[english,french,review]{cjs-rcs-article}
  %% Commandes réservées aux éditeurs de La revue. Laisser en
  %% commentaires sans les supprimer.
  % \jvol{00}
  % \jissue{00}
  % \jyear{YYYY}
  % \firstpage{1}
  % \received{YYYY-MM-DD}
  % \accepted{YYYY-MM-DD}
  % \licence{}

  %% Chargez ici les éventuels paquetages additionnels. Utilisez
  %% uniquement des paquetages absolument requis pour le manuscript.
  % \usepackage{}

  %% Informations de la page de titre et des parties annexes.
  %% Remplacez le texte entre < >. Supprimez les caractères < >.
  %%
  %% TITRE DE L'ARTICLE. Utiliser un titre concis et descriptif; il ne
  %% doit pas excéder deux lignes. Insérez les changements de ligne
  %% manuellement avec \\ (avant les déterminants et prépositions). Si
  %% le titre est trop long pour l'entête, fournissez une version
  %% courte en argument optionnel (entre les crochets [ ]).
  \title{<Titre de l'article>}
  % \title[Version courte du titre]{Ceci est un exemple de long titre \\
  %   avec saut de ligne manuel}

  %% AUTEURS ET AFFILIATIONS. L'identification des auteurs se fait
  %% par des paires de commandes \author et \affil, une par auteur.
  %% Saisissez le nom de l'auteur entièrement en minuscules (sauf pour
  %% les premières lettres et les initiales). Placez le nom (de
  %% famille) dans une commande \surname. L'argument optionnel permet
  %% d'indiquer l'identifiant ORCID de l'auteur, son adresse courriel
  %% et son statut d'auteur correspondant.
  \author[orcid=<id>, email=<courriel>, corresponding]
         {<Prénom> \surname{<Nom>}}
  \affil{<Département, Université, Ville, Pays>}
  % \author[orcid=<id>, email=<courriel>]
  %        {<Prénom> \surname{<Nom>}}
  % \affil{<Département, Université, Ville, Pays>}
  % \author[orcid=<id>, email=<courriel>]
  %        {<Prénom> \surname{<Nom>}}
  % \affil{<Département, Université, Ville, Pays>}

  %% (Optionnel - décommenter pour utiliser) Si la liste des auteurs
  %% est trop longue pour l'entête de l'article, utilisez la commande
  %% \runningauthor pour fournir une identification plus courte de la
  %% forme «Nom et collab.»
  % \runningauthor{<Nom> et collab.}

  %% RÉSUMÉS. Vous devez fournir un résumé en français et un résumé en
  %% anglais. L'ordre est sans importance; ils seront positionnés sur
  %% la page de titre automatiquement.
  \begin{englishabstract}
    <Insert your English abstract here. The abstract should be concise
    (up to 20 lines long) and descriptive. Avoid symbols as much as
    possible. Formulas are strongly discouraged, and citations should
    be avoided.>
    [Nous fournirons un résumé en anglais pour un article en français
    aux auteurs qui ne sont pas en mesure de le rédiger eux-mêmes.]
  \end{englishabstract}

  \begin{frenchabstract}
    <Insérez le texte du résumé en français ici. Le résumé doit être
    concis (20 lignes ou moins) et descriptif. Évitez autant que
    possible les symboles. Les équations mathématiques sont fortement
    déconseillées et les citations sont à éviter.>
  \end{frenchabstract}

  %% MOTS-CLÉS. Fournissez une liste de mots-clés en ordre
  %% alphabétique. Seul le premier mot-clé doit débuter par une
  %% majuscule. N'utilisez aucune ponctuation.
  \begin{keywords}
  \item <Mot-clé 1>
  \item <mot-clé 2>
  \item <mot-clé 3>
  \end{keywords}

  %% CLASSIFICATION MATHÉMATIQUE. Fournissez la classification
  %% mathématique par matières de l'article selon le classement
  %% MSC2020 (https://msc2020.org). Celle-ci est habituellement
  %% divisée en niveaux primaire et secondaire. Séparez les matières
  %% par des virgules, sans aucune autre ponctuation.
  \begin{classification}
  \item[Primaire] <rubrique 1>, <rubrique 2>
  \item[Secondaire] <rubrique 1>, <rubrique 2>
  \end{classification}

  %% MATÉRIEL ADDITIONNEL (OBLIGATOIRE si l'article est accompagné de
  %% matériel additionnel - commenter si inutilisé). Fournissez la
  %% liste du matériel additionnel disponible en ligne. Ce matériel
  %% doit être cité dans l'article. Ce contenu est affiché par la
  %% commande \makebackmatter.
  \begin{supplement}
    <Décrivez brièvement le matériel additionnel.>
  \end{supplement}

  %% PARTAGE DES DONNÉES (OBLIGATOIRE). Conformément à la politique
  %% éditoriale «Expect Data Sharing» de Wiley, vous devez indiquer
  %% les conditions de partage des données de recherche, ou l'absence
  %% de celles-ci, le cas échéant. Fournir un lien permanent, comme un
  %% identificateur d'objets numériques (DOI) ou une adresse URL. Ce
  %% contenu est affiché par la commande \makebackmatter.
  \begin{sharing}
    <Précisez les conditions de partage des données de recherche, ou
    encore leur absence, le cas échéant.>
  \end{sharing}

  %% REMERCIEMENTS (optionnel - commenter si inutilisé). Inscrivez vos
  %% remerciements ici. Attention de ne pas répéter la rubrique de
  %% financement de la recherche, ci-dessous. Ce contenu est affiché
  %% par la commande \makebackmatter.
  \begin{acknowledgements}
    <Insérez vos remerciements ici.>
  \end{acknowledgements}

  %% FINANCEMENT DE LA RECHERCHE (optionnel - commenter si inutilisé).
  %% Inscrivez les informations sur le financement de la recherche
  %% sous forme de texte libre. Indiquez le nom des agences de
  %% financement au long. Ce contenu est affiché par la commande
  %% \makebackmatter.
  \begin{funding}
    <Inscrivez les informations sur le financement ici.>
  \end{funding}

\begin{document}

%% Page de titre (OBLIGATOIRE - ne supprimez pas la ligne suivante!)
\maketitle                      % composer la page de titre

%% Introduction numérotée (OBLIGATOIRE)
\section{Introduction}
\label{sec:introduction}        % exemple d'étiquette de section

<Débutez l'article par une introduction.>

%% Sections additionnelles
\section{<Titre de section>}

<Utilisez les commandes LaTeX pour diviser en sections, sous-sections,
sous-sous-sections, etc. Écrivez les titres avec uniquement une
majuscule initiale.>

%% >>>>>>>>>>>>>>>>>>>>>>>>>>>>>>>>>>>>>>>>>>>>>>>>>>>>>>>>>>>>>>>>>>
%% Ce bloc contient des exemples de déclarations pour diverses
%% structures fréquemment employées dans des articles. Utilisez-les
%% comme bon vous semble et supprimez ce qui est inutile.
\section{Exemples de déclarations}

%% Références croisées.
Voici deux exemples de références croisées: d'abord à la
\autoref{sec:introduction}, puis à l'équation~\eqref{eq:continuity}.

%% Environnements de type théorème. Consultez aussi la documentation
%% du paquetage amsthm.
\begin{definition}
  \label{def:cdf}
  La fonction de répartition $F$ d'une variable aléatoire réelle $X$
  est définie, pour tout $x \in \Rset$, comme $F(x) = \Pr(X \leq x)$.
\end{definition}

\begin{theorem}
  \label{thm:continuity}
  Soit $F$ la fonction de répartition d'une variable aléatoire réelle
  tel qu'à la \autoref{def:cdf}. Alors $F$ est non décroissante,
  continue par la droite et elle satisfait
  \begin{equation}
    \label{eq:continuity}
    \lim_{x \to -\infty} F(x) = 0, \quad \lim_{x \to \infty} F(x) = 1.
  \end{equation}
\end{theorem}

\begin{lemma}
  \label{lem:continuity}
  Un lemme qui permet de démontrer le \autoref{thm:continuity}.
\end{lemma}

\begin{proof}
  La démonstration du \autoref{lem:continuity}.
\end{proof}

\begin{proof}[Démonstration de la continuité de $F$]
  Démonstration du \autoref{thm:continuity}.
\end{proof}

\begin{corollary}
  Un corollaire important du \autoref{thm:continuity}.
\end{corollary}

\begin{proposition}
  \label{prop:foobar}
  Un résultat de moindre importance peut être présenté sous forme de
  proposition.
\end{proposition}

\begin{algorithm}[Algorithme d'Euclide]
  \label{algo:euclid}
  Calculer le plus grand commun diviseur de deux entiers positifs $m$
  et $n$ avec, sans perte de généralité, $n < m$.
  \begin{enumerate}
  \item Diviser $m$ par $n$ et poser $r$ égal au reste de la division.
    (Nous avons alors $0 \leq r < n$.)
  \item Si $r = 0$, retourner la valeur $n$.
  \item Poser $m \leftarrow n$, $n \leftarrow r$, puis retourner à
    l'étape 1. \qed
  \end{enumerate}
\end{algorithm}

\begin{remark}
  Cette version de \autoref{algo:euclid} est adaptée de Knuth, D. E.,
  \emph{The Art of Computer Programming}, vol.~1, 1997.
\end{remark}

%% Figure. L'image est tirée du paquetage mwe; elle devrait être
%% disponible sur votre système.
\begin{figure}
  \centering
  \includegraphics[width=0.4\linewidth]{example-grid-100x100pt}
  \caption{Un exemple de figure}
  \label{fig:example-grid}
\end{figure}

Allons-y d'une référence à la \autoref{fig:example-grid}.

%% Tableau. La commande \np (raccourci de \numprint) du paquetage
%% numprint ajoute automatiquement le séparateur approprié selon la
%% langue à tous les groupes de trois chiffres. Le séparateur décimal
%% en français est la virgule. Si vous devez inscrire des notes
%% de contenu dans un tableau, utilisez le paquetage
%% threeparttable (tel que recommandé dans l'entrée de la Foire
%% aux questions de TeX https://texfaq.org/FAQ-footintab).
\begin{table}
  \caption{Liste des dix pays ayant la plus grande publication
    intérieure brute (\textsc{gnp}) de recherche en statistique. Le
    classement est effectué selon la variable \textsc{pag}$^{\star}$.}
  \label{tab:gnp}
  \centering
  \begin{tabular}{rlrrrrrr}
    \toprule
    Rang & Pays & \textsc{pag}$^{\star}$ & \textsc{pag}
    & \textsc{art}$^{\star}$ & \textsc{art} \\
    \midrule
    1  & États-Unis     & \np{109338} & \np{60369} & \np{7240} & \np{4061} \\
    2  & Royaume-Uni    & \np{ 12597} & \np{ 7504} & \np{ 884} & \np{ 538} \\
    4  & Australie      & \np{  7872} & \np{ 4261} & \np{ 578} & \np{ 323} \\
    5  & Allemagne      & \np{  6782} & \np{ 4500} & \np{ 456} & \np{ 306} \\
    7  & Japon          & \np{  2865} & \np{ 1880} & \np{ 241} & \np{ 163} \\
    8  & Pays-Bas       & \np{  2864} & \np{ 1702} & \np{ 191} & \np{ 116} \\
    9  & Inde           & \np{  2559} & \np{ 1395} & \np{ 275} & \np{ 151} \\
    10 & Israël         & \np{  2097} & \np{ 1160} & \np{ 148} & \np{  83} \\
    \bottomrule
  \end{tabular}
\end{table}

Voici une référence au \autoref{tab:gnp}.
%% <<<<<<<<<<<<<<<<<<<<<<<<<<<<<<<<<<<<<<<<<<<<<<<<<<<<<<<<<<<<<<<<<<

%% Conclusion (optional - comment out if unused)
\section{Conclusion}

<Vous pouvez terminer votre article avec une conclusion dans la mesure
où celle-ci propose quelque chose de nouveau. Il est inutile d'inclure
une conclusion qui ne fait que répéter le contenu de l'article.>

%% Parties annexes (OBLIGATOIRE - ne supprimez pas la ligne suivante!)
\makebackmatter                 % composer les parties annexes

%% Annexes (optionnel - commenter si inutilisé)
\appendix                       % début des annexes

\section{<Titre de l'annexe>}

<Les annexes sont dévolues aux démontrations de théorèmes ou autres
longs développements mathématiques. Vous devez citer les annexes dans
l'article.>

%% Liste des références. Citez avec les commandes du paquetage natbib.
%% Insérez les noms des bases de données bibliographiques (fichiers
%% .bib) dans la commande \bibliography, sans extension et séparés
%% par des virgules.
\bibliography{}

\end{document}
